\documentclass[11pt]{article}
\usepackage{icmlko}
\begin{document}

\title{상태 층을 두 달 만에 열었더니 기억이 열이레였다 --- 설계가 열 배 길게 잡혀 있었다}
\author{월드모델 실험실}
\date{노트 620 · 621 · 622 · 623 · 2026년 8월 5일}
\maketitle

\begin{abstract}
이 저장소의 로드맵은 블록 $3$ 을 ``엔티티 메모리 스트림 --- 브랜드/공간
단위 관찰$+$성찰 노드, \texttt{score = recency(0.995\^{}t) + importance +
relevance}''로 두 달 전에 적어 두었다. \textbf{한 줄도 안 지어졌고 그
이유가 기록돼 있었다 --- ``상태 층 $0$ 행''.}

왜 $0$ 행인지 재 보니 팝업 레코드 쪽은 원리상 못 잰다 --- 채점 풀의 IP
$68$ 개 중 \textbf{$2$ 회 이상이 $7$ 개}다. 그런데 \textbf{다른 데에
있었다} --- \texttt{popga} 이벤트의 \texttt{spot\_id} $\times$ 월이
\textbf{$63{,}826$ 관측 $\cdot$ $7{,}701$ 스팟}이고 $6$ 개월 이상이
$4{,}959$ 스팟이다.

\begin{center}
\begin{tabular}{lrrl}
\hline
마지막 관측과의 간격 & $n$ & MAE 이득 & 스팟 군집 $95\%$ \\
\hline
\textbf{$1$ 달} & $47{,}638$ & $\mathbf{0.0724}$ & $[0.068, 0.077]$ \\
$2$ 달 & $5{,}118$ & $0.0147$ & $[0.010, 0.019]$ \\
$3$ 달 & $1{,}571$ & $0.0105$ & $[0.001, 0.020]$ \\
$4\!\sim\!6$ 달 & $238\!\sim\!744$ & --- & \textbf{못 잰다} \\
$7\!\sim\!9$ 달 & $289$ & $0.0000$ & \textbf{진다} \\
\hline
\end{tabular}
\end{center}

\textbf{한 달 지나면 이득의 $80\%$ 가 사라진다} --- 첫 구간 반감기
\textbf{$13$ 일}. 로드맵이 적어 둔 $0.995^t$ 은 반감기 \textbf{$139$ 일}이다.
\end{abstract}

\begin{figure}[t]
\centering
\includegraphics[width=\textwidth]{figs/f1.pdf}
\caption{\textbf{가} --- 초록만 잴 수 있다(세 칸).
\textbf{나} --- 제일 가까울 때도 반만 믿는다.
\textbf{다} --- 설계가 열 배 길게 잡혀 있었다.}
\end{figure}

\section{엔티티 고정효과는 없다}

먼저 갈랐다 --- $y(s,t) = m(t) + a(s) + e(s,t)$ 로 놓고 시장 수준 $m$ 과
스팟 고유 오프셋 $a$ 를 따로 만들었다.

\begin{center}
\begin{tabular}{lr}
\hline
팔 & MAE \\
\hline
$m(t)$ 시장 수준만 & $\mathbf{0.3513}$ \\
$m(t)+a(s)$ 시장 $+$ 스팟 고유 & $0.4833$ \\
$m(t{-}1)+a(s)$ (배포 가능) & $0.4953$ \\
$m(t{-}1)+$ 축소한 $a(s)$ & $0.3654$ \\
\hline
\end{tabular}
\end{center}

\textbf{고유 오프셋을 더하면 크게 나빠진다}($+0.132$, 군집
$[+0.126,+0.138]$). 그리고 그 $a$ 를 $0$ 쪽으로 \emph{축소}하면 손해가
거의 다 회복된다 --- \textbf{$a$ 는 정보가 아니라 잡음이다.}

\begin{quote}
\textbf{스팟에는 안정적인 ``수준''이 없다.} 팝업은 열리고 치솟고 닫힌다 ---
생애 평균 편차는 이번 달과 아무 상관이 없다. \textbf{메모리 스트림이
전제하는 ``엔티티가 지속하는 상태를 갖는다''가 이 자료에서 거짓이다.}
\end{quote}

\section{그런데 마지막 관측 하나는 값이 있다}

연속 세 달이 있는 지점($n{=}13{,}703$ 채점 $\cdot$ 스팟 $4{,}511$)에서

\begin{center}
\begin{tabular}{lrr}
\hline
 & MAE & 앞 칸 대비 \\
\hline
$m(t{-}1)$ 시장만 & $0.3236$ & --- \\
\textbf{$+$ lag$1$} & $\mathbf{0.2637}$ & $\mathbf{-0.0599}$ \\
$+$ lag$2$ & $0.2607$ & $-0.0030$ \\
$+$ lag$3$ & $0.2601$ & $-0.0005$ \\
\hline
\end{tabular}
\end{center}

\textbf{lag$2 \cdot$ lag$3$ 도 유의하지만 이득이 $20$ 배 $\cdot$ $120$ 배
작다.} 능형 계수도 같은 말을 한다 --- lag$1$ $\mathbf{0.607}$ 대 lag$2$
$0.032$ 대 lag$3$ $0.018$. \textbf{사실상 $1$ 차다.}

\section{앞의 둘이 왜 갈렸나 --- 그것이 감쇠 곡선이었다}

전체 표본($n{=}56{,}125$, 간격 있는 관측 포함)에서는 ``직전 관측''이 시장
수준을 \textbf{못 이겼다}($+0.003$, $0$ 을 문다). 연속 세 달에서는
$-0.0599$ 로 크게 이겼다. \textbf{모순이 아니라 표본이 달랐다} --- 앞은
간격이 섞여 있고 뒤는 간격이 $1$ 로 고정이다.

간격별로 다시 재니 위 표가 나왔다. \textbf{$g{=}1$ 에서 $g{=}2$ 로 이득이
$0.0724 \to 0.0147$ 로 $80\%$ 사라진다.} 지수로 읽으면 반감기
$\mathbf{13.2}$ 일이다.

\begin{quote}
\textbf{``상태''가 아니라 ``최근 관측''이다.} 엔티티가 들고 다니는 잠재
수준은 없고, 어제 본 값이 오늘을 조금 말해 줄 뿐이며 그 조금이 이 주면
반으로 준다.
\end{quote}

\section{제일 가까울 때도 반만 믿는다}

혼합 $\alpha \cdot y(s,t{-}1) + (1-\alpha) \cdot m(t{-}1)$ 의 최적
$\alpha$ 는 $g{=}1$ 에서 \textbf{$0.50$} 이다(연속 세 달 표본에서는
$0.625$). \textbf{제일 좋은 조건에서도 자기 과거의 몫이 절반}이고 나머지
절반은 시장 수준이다.

\begin{quote}
\textbf{잴 수 있는 칸은 셋뿐이다}($g{=}1,2,3$). $g{=}4\!\sim\!6$ 은
$n$ 이 $238\!\sim\!744$ 라 군집 구간이 $0$ 을 문다. 그 칸들에서 $\alpha$ 가
$0.40 \to 0.20 \to 0.40$ 으로 튀는 것은 \textbf{순서가 아니라 잡음}이다
(노트 $608$: 작은 칸이 구간 끝을 정한다). \textbf{잴 수 있는 세 칸은
단조 감소한다.}
\end{quote}

\section*{조항}

\begin{quote}
\textbf{설계를 적기 전에 상수를 잰다.} 로드맵은 두 달 전에
\texttt{recency(0.995\^{}t)} 를 적어 두었다 --- 근거 없이 그럴듯해서 고른
숫자였고 \textbf{열 배 길다}. 그 상수 위에 관찰$\cdot$성찰 노드를
지었으면 \textbf{$5$ 개월 전 관찰이 지금과 같은 무게}로 들어갔을 것이다.
\textbf{$0$ 행이었던 것이 다행이었다.}
\end{quote}

\section*{남은 것}

\begin{itemize}
\item \textbf{이것은 방문 트래픽 상태이지 팝업 성과 상태가 아니다.}
\texttt{spot\_id} $\to$ 팝업 레코드 매핑을 안 했다. \textbf{$13$ 일이
성과 쪽에도 같은지는 안 쟀다.}
\item \textbf{브랜드 쪽은 아직 $0$ 이다.} 공간(\texttt{spot\_id})만 열렸다.
로드맵 블록 $3$ 은 브랜드와 공간 둘 다인데 브랜드 스트림은 여전히 조인
키가 없다(레코드의 \texttt{entities} 가 전부 \texttt{unresolved:} 다).
\item \textbf{$m(t)$ 를 쓰는 팔은 배포 불가다.} 그 달 전역 중앙값은
그 달이 끝나야 안다. 배포 기준선은 $m(t{-}1)$ 이고 그 차가
$0.3513$ 대 $0.3638$ 이다 --- \textbf{시장 수준 예측 자체가 별도 과제}다.
\end{itemize}

\end{document}
